\SourceLeaf{388}%
regard much of the contents only as dissertations in which elementary subjects are treated with more breadth than depth. It would, however, be inappropriate to subject these works. which belong to a time long past, to detailed criticism now.

The new continuation, by contrast, appears far more substantial, and military readers, historians and geographers will find much to learn from it. We mention here only the account of the Battle of Marengo, the history of the campaign in Germany in 1800 (which fills almost the whole fifth volume), and the military description of the Danube basin, all illustrated with a great many maps and plans. There are also the proceedings of a commission specially appointed to consider the most suitable way of representing terrain. They examine the subject from many angles and illustrate it with a considerable number of sample drawings made by different methods. Readers will find it interesting to read, in the correspondence of Count \textsc{de Gisors} with his father, the Duke \textsc{de Belle-Isle}, the conversations that the former had with \textsc{Frederick} the Second on military matters a year before the outbreak of the Seven Years' War. The same goes for a series of previously unpublished letters of \textsc{Louis} XIV, some in his own hand, which are also accompanied by a facsimile. Readers may, however, differ in how far they can see the letters from the point of view the editors wish their French readers to adopt when they remark in the introduction, \textsc{Montesquieu} \textit{regarde comme le devoir de tout écrivain homme de bien de contribuer, autant qu’il est en lui, à donner à ses concitoyens des raisons d’aimer ceux à qui ils doivent obéir. Rien n’entre mieux dans cette noble pensée de} \textsc{Montesquieu}, \textit{que la publication de ces lettres et de celles qui pourront les suivre, et tout le monde reconnaîtra dans les successeurs du grand roi tout ce que son coeur avait de paternel et son âme d’héroïque.} Finally, we must also mention the reports in the 3rd and 4th volumes on the new large map of France, whose production was ordered by a royal ordinance of 6 August 1817. The original plan was to carry out the survey at a scale of 1 to 10000 and the engraving at a scale of 1 to 50000. The total number of sheets was estimated at 611, each 800 millimetres wide and 500 millimetres high, and it was thought that the whole work could be completed in 20 years. This plan was soon abandoned, however, and the scales were reduced for the survey
